\subsubsection{纤维多重度的矩阵积表示与张量化 MOM 编译：有限窗逆码下的非交换计数}\label{subsubsec:pom-multiplicity-matrix-product-tensor-mom}
在正文的碰撞谱口径中，\(S_q(m)=\sum_{x\in X_m}d_m(x)^q\) 被视为一个“\(q\)-重纤维积”计数对象（定义 \ref{def:pom-moment-q}），并通过有限状态 moment-kernel 统一纳入 Perron--Frobenius/\(\zeta\) 接口（命题 \ref{prop:pom-moment-kernel-exists} 及定义 \ref{def:pom-moment-pressure}）。
本段补充一个更细的结构层：在存在\emph{有限窗逆码}（finite-window inverse code）的制度下，单个纤维体积 \(d_m(x)\) 本身就具有矩阵积计数式；并且对任意整数阶 \(q\)，\((\mathrm{MOM}_q)\) 的碰撞核可显式写成“语法两态骨架 \(\times\) 逆码转移的 \(q\)-重同步直积”的 Kronecker 耦合。

\paragraph{有限窗逆码 \texorpdfstring{$\Rightarrow$}{=>} 纤维体积的矩阵积计数}
\begin{theorem}[纤维多重度的矩阵积表示与最小歧义态维数]\label{thm:pom-fiber-multiplicity-matrix-product}
\leanverified{paper\_pom\_fiber\_multiplicity\_matrix\_product}
设 \(\Fold_\infty:\Omega_\infty=\{0,1\}^{\NN}\to X_\infty\) 为二元折叠因子，且 \(X_\infty\subseteq\{0,1\}^{\NN}\) 为黄金均值移位（禁止相邻 \(11\)）。
假设 \(\Fold_\infty\) 存在一个有限窗逆码 \(\Psi\)，使得“给定输出前缀 \(x_1\cdots x_m\)”时其全部合法预像可由一个有限状态确定化系统逐位更新并计数（等价地，可取一个 follower-separated 的右解析表示并在其上做计数读出）。
则存在有限状态集合 \(\mathcal S\)、边界向量 \(u,v\in\NN^{\mathcal S}\)，以及两矩阵
\[
T_0,T_1\in M_{|\mathcal S|}(\NN),
\]
使得对任意合法宏观词 \(x=x_1\cdots x_m\in X_m\)（\(m\ge 1\)）有精确计数式
\[
\boxed{\;
d_m(x)=u^{\mathsf T}T_{x_1}T_{x_2}\cdots T_{x_m}v\; }.
\]
此外，在所有满足上述表示的选择中，最小可能维数 \(|\mathcal S|\) 等于右歧义类（follower ambiguity classes）的有限完备集合基数；
因此当 \(\Psi\) 取为由最小在线延迟/最小同步结构诱导的规范逆码时，该最小维数由 \(\Fold_\infty\) 唯一确定。
\end{theorem}
\begin{proof}
由假设，逆码可实现为一个确定性有限状态计数机：其状态记录“在给定输出前缀下仍可能的歧义类”，读入下一个输出符号 \(b\in\{0,1\}\) 后作确定更新。
令 \(\mathcal S\) 为该计数机的（可达）状态集，并令 \((T_b)_{b\in\{0,1\}}\) 为按输出符号分解的转移计数矩阵（每一步的计数语义由逆码产生的分支数给出，因而条目为非负整数）。
取 \(u\) 为初始状态分布（由边界/起始允许态确定），取 \(v\) 为末端读出向量（对允许终止态求和），则长度 \(m\) 的输出词 \(x_1\cdots x_m\) 的全部预像数正是沿该确定性计数机的路径计数，从而得到矩阵积公式。
\par
最小性结论是确定性有限自动机的标准最小化：将状态按“拥有相同 follower 集（同一切右延拓族）”的等价关系合并得到最小右解析 follower-separated 表示，其等价类数即为所述右歧义类完备集合基数；任何矩阵积表示必须区分这些等价类，故其维数不小于该基数，而由最小化构造可达该下界。
\end{proof}

\paragraph{张量化 MOM 编译：黄金语法骨架 \texorpdfstring{$\otimes$}{x} \(q\)-重逆码同步直积}
\begin{proposition}[MOM\texorpdfstring{$\,_q$}{\_q} 的张量分解型碰撞核]\label{prop:pom-momq-tensor-kernel-construction}
\leanverified{paper\_pom\_momq\_tensor\_kernel\_construction}
在定理 \ref{thm:pom-fiber-multiplicity-matrix-product} 的设定下，令
\[
G_0:=\begin{pmatrix}1&0\\[2pt]1&0\end{pmatrix},
\qquad
G_1:=\begin{pmatrix}0&1\\[2pt]0&0\end{pmatrix}
\in M_2(\{0,1\})
\]
为黄金均值两态骨架的符号分解转移矩阵（满足 \(G_0+G_1=\bigl(\begin{smallmatrix}1&1\\1&0\end{smallmatrix}\bigr)\)）。
对任意整数 \(q\ge 1\)，定义显式的 \(q\)-阶碰撞转移核
\[
A_q^{\mathrm{ten}}
\;:=\;
G_0\otimes T_0^{\otimes q}\;+\;G_1\otimes T_1^{\otimes q}
\ \in\ M_{2|\mathcal S|^q}(\NN).
\]
令 \(e_0:=(1,0)^{\mathsf T}\in\NN^2\)、\(\mathbf 1:=(1,1)^{\mathsf T}\in\NN^2\)，并取边界向量
\[
\alpha_q:=e_0\otimes u^{\otimes q},
\qquad
\beta_q:=\mathbf 1\otimes v^{\otimes q}.
\]
则对所有 \(m\ge 1\) 有线性递推计数表示
\[
\boxed{\;
S_q(m)=\sum_{x\in X_m}d_m(x)^q
\;=\;
\alpha_q^{\mathsf T}\bigl(A_q^{\mathrm{ten}}\bigr)^{m}\beta_q\; }.
\]
因此 \((\mathrm{MOM}_q)\) 的”有限状态编译”可取为张量分解型结构：\(q\)-体重合约束在核层面等同于逆码转移的 \(q\)-重同步直积，再与黄金语法两态骨架作 Kronecker 耦合。
\leanverified{collision\_kernel\_dimensions}
\end{proposition}
\begin{proof}
由定理 \ref{thm:pom-fiber-multiplicity-matrix-product}，
\(
d_m(x)=u^{\mathsf T}T_{x_1}\cdots T_{x_m}v
\)。
于是
\[
d_m(x)^q
=
\bigl(u^{\otimes q}\bigr)^{\mathsf T}\,
\bigl(T_{x_1}^{\otimes q}\bigr)\cdots
\bigl(T_{x_m}^{\otimes q}\bigr)\,
v^{\otimes q}.
\]
对所有合法词 \(x\in X_m\) 求和可由黄金均值两态骨架的路径展开实现：每一步选择符号 \(b\in\{0,1\}\) 并按 \(G_b\) 更新语法态，同时在 \(\mathcal S^{\times q}\) 上施加 \(T_b^{\otimes q}\) 的同步更新。
因此加权路径总和恰为 \(\alpha_q^{\mathsf T}(A_q^{\mathrm{ten}})^m\beta_q\)。证毕。
\end{proof}

\begin{proposition}[MOM 张量核的 Schur 补降维行列式公式]\label{prop:pom-momq-tensor-kernel-schur-determinant}
\leanverified{paper\_pom\_momq\_tensor\_kernel\_schur\_determinant}
在命题 \ref{prop:pom-momq-tensor-kernel-construction} 的设定下，对任意 \(q\ge 1\) 与任意 \(z\in\CC\) 都有恒等式
\[
\boxed{\;
\det(I-zA_q^{\mathrm{ten}})
=\det\!\Bigl(I-zT_0^{\otimes q}-z^2\,(T_1T_0)^{\otimes q}\Bigr)\; }.
\]
因此对应的 \(\zeta\)-函数满足
\[
\zeta_{A_q^{\mathrm{ten}}}(z)
:=\det(I-zA_q^{\mathrm{ten}})^{-1}
=\det\!\Bigl(I-zT_0^{\otimes q}-z^2\,(T_1T_0)^{\otimes q}\Bigr)^{-1}.
\]
\end{proposition}
\begin{proof}
由 \(G_0,G_1\) 的显式形状直接展开得
\[
A_q^{\mathrm{ten}}
=
\begin{pmatrix}
T_0^{\otimes q} & T_1^{\otimes q}\\[2pt]
T_0^{\otimes q} & 0
\end{pmatrix}.
\]
从而
\[
I-zA_q^{\mathrm{ten}}
=
\begin{pmatrix}
I-zT_0^{\otimes q} & -zT_1^{\otimes q}\\[2pt]
-zT_0^{\otimes q} & I
\end{pmatrix}.
\]
对右下块 \(I\) 作 Schur 补，得
\[
\det(I-zA_q^{\mathrm{ten}})
=\det\!\Bigl(I-zT_0^{\otimes q}-z^2\,T_1^{\otimes q}T_0^{\otimes q}\Bigr).
\]
最后由 Kronecker 幂的乘法性 \(T_1^{\otimes q}T_0^{\otimes q}=(T_1T_0)^{\otimes q}\) 得结论。证毕。
\end{proof}

\begin{proposition}[非交换 Fibonacci 多项式给出的 \texorpdfstring{$(A_q^{\mathrm{ten}})^m$}{(Aq^ten)^m} 显式闭式]\label{prop:pom-momq-tensor-kernel-noncommutative-fibonacci-powers}
沿用上述记号。定义矩阵序列 \((C_m^{(q)})_{m\ge 0}\subseteq M_{|\mathcal S|^q}(\NN)\) 为
\[
C_0^{(q)}:=I,\qquad
C_1^{(q)}:=T_0^{\otimes q},\qquad
C_m^{(q)}:=T_0^{\otimes q}C_{m-1}^{(q)}+(T_1T_0)^{\otimes q}C_{m-2}^{(q)}\quad(m\ge 2),
\]
并约定 \(C_{-1}^{(q)}=C_{-2}^{(q)}=0\)。
则对任意 \(m\ge 0\) 都有严格恒等式
\[
\boxed{\;
\bigl(A_q^{\mathrm{ten}}\bigr)^m
=
\begin{pmatrix}
C_m^{(q)} & C_{m-1}^{(q)}\,T_1^{\otimes q}\\[2pt]
T_0^{\otimes q}\,C_{m-1}^{(q)} & \mathbf 1_{\{m=0\}}I+T_0^{\otimes q}\,C_{m-2}^{(q)}\,T_1^{\otimes q}
\end{pmatrix}\; }.
\]
\leanverified{paper\_pom\_momq\_tensor\_kernel\_noncommutative\_fibonacci\_powers}
\end{proposition}
\begin{proof}
当 \(m=0,1\) 由定义直接核对成立。
设 \(m\ge 1\) 且结论对该 \(m\) 成立。由
\(
A_q^{\mathrm{ten}}=\bigl(\begin{smallmatrix}T_0^{\otimes q}&T_1^{\otimes q}\\ T_0^{\otimes q}&0\end{smallmatrix}\bigr)
\)
与块乘法得
\[
\bigl(A_q^{\mathrm{ten}}\bigr)^{m+1}
=
\begin{pmatrix}
T_0^{\otimes q}C_m^{(q)}+(T_1T_0)^{\otimes q}C_{m-1}^{(q)} & C_m^{(q)}T_1^{\otimes q}\\[2pt]
T_0^{\otimes q}C_m^{(q)} & T_0^{\otimes q}C_{m-1}^{(q)}T_1^{\otimes q}
\end{pmatrix}.
\]
其中左上块恰由递推给出 \(C_{m+1}^{(q)}\)，其余三块与所示公式在 \(m+1\) 时一致。
归纳完成。证毕。
\end{proof}

\paragraph{张量折叠的递推结果式、Hadamard 极点与 Galois 刚性}
在命题 \ref{prop:pom-momq-tensor-kernel-construction} 给出的 Kronecker 型碰撞核之上，若把若干条折叠系统作外部张量并要求矩谱按尺度逐点相乘，则最小递推特征式不再是抽象存在量，而受一个显式结式严格控制；在无碰撞制度下，这一控制进一步提升为精确等式、多因子 Hankel 秩严格乘法、Hadamard 极点的完全可见性，以及产品谱根的 Galois 次数刚性。

\begin{theorem}[张量折叠碰撞矩的结果式递推律]\label{thm:pom-tensor-fold-resultant-recursion}
固定整数 \(q\ge 1\)。
设两条折叠系统在矩阶 \(q\) 上的碰撞矩序列分别为
\[
\bigl(S_q^{(1)}(m)\bigr)_{m\ge 0},
\qquad
\bigl(S_q^{(2)}(m)\bigr)_{m\ge 0},
\]
并假设其张量折叠满足严格乘法律
\[
S_q^{(\otimes)}(m)=S_q^{(1)}(m)\,S_q^{(2)}(m)\qquad(m\ge 0).
\]
记 \(P_q^{(i)}(\lambda)\in\ZZ[\lambda]\) 为 \(\bigl(S_q^{(i)}(m)\bigr)_{m\ge 0}\) 的首一最小递推特征多项式，并令 \(d_i:=\deg P_q^{(i)}\)（\(i=1,2\)）。
定义
\[
\widehat P_q^{(1)}(\lambda,u):=u^{d_1}P_q^{(1)}(\lambda/u)\in\ZZ[\lambda,u],
\qquad
R_q(\lambda):=\operatorname{Res}_{u}\!\bigl(\widehat P_q^{(1)}(\lambda,u),\,P_q^{(2)}(u)\bigr)\in\ZZ[\lambda].
\]
则张量折叠序列 \(\bigl(S_q^{(\otimes)}(m)\bigr)_{m\ge 0}\) 必满足某条整系数线性递推；
若其首一最小递推特征多项式记为 \(P_q^{(\otimes)}(\lambda)\)，则
\[
P_q^{(\otimes)}(\lambda)\mid R_q(\lambda)\qquad\text{于 }\QQ[\lambda].
\]
此外，若在 \(\overline{\QQ}\) 中把
\[
P_q^{(1)}(\lambda)=\prod_{i=1}^{d_1}(\lambda-\alpha_i),
\qquad
P_q^{(2)}(\lambda)=\prod_{j=1}^{d_2}(\lambda-\beta_j)
\]
按重数分解，则有严格因式分解
\[
R_q(\lambda)=\prod_{i=1}^{d_1}\prod_{j=1}^{d_2}(\lambda-\alpha_i\beta_j).
\]
\end{theorem}
\leanverified{paper\_pom\_tensor\_fold\_resultant\_recursion}
\begin{proof}
取 \(P_q^{(i)}\) 的伴随矩阵 \(C_i\in M_{d_i}(\ZZ)\)。
由于 \(P_q^{(i)}\) 首一且为最小递推特征多项式，存在向量 \(u_i,v_i\) 使
\[
S_q^{(i)}(m)=u_i^{\mathsf T}C_i^{\,m}v_i\qquad(m\ge 0),
\qquad
\chi_{C_i}(\lambda)=P_q^{(i)}(\lambda).
\]
由张量折叠的乘法律，
\[
S_q^{(\otimes)}(m)
=\bigl(u_1\otimes u_2\bigr)^{\mathsf T}
\bigl(C_1\otimes C_2\bigr)^{m}
\bigl(v_1\otimes v_2\bigr).
\]
因此 \(\bigl(S_q^{(\otimes)}(m)\bigr)\) 是矩阵 \(C_1\otimes C_2\) 的线性读出；于是任一满足
\(
F(C_1\otimes C_2)=0
\)
的多项式 \(F\in\QQ[\lambda]\) 都诱导一条递推
\(
F(E)S_q^{(\otimes)}\equiv 0
\)。

在 \(\overline{\QQ}\) 中，\(C_1\) 与 \(C_2\) 的特征值多重集分别为 \(\{\alpha_i\}_{i=1}^{d_1}\) 与 \(\{\beta_j\}_{j=1}^{d_2}\)，故
\[
\chi_{C_1\otimes C_2}(\lambda)
=\prod_{i=1}^{d_1}\prod_{j=1}^{d_2}(\lambda-\alpha_i\beta_j).
\]
另一方面，
\[
\widehat P_q^{(1)}(\lambda,u)
=u^{d_1}P_q^{(1)}(\lambda/u)
=\prod_{i=1}^{d_1}(\lambda-\alpha_i u),
\]
因而结式恒等式给出
\[
R_q(\lambda)
=\operatorname{Res}_{u}\!\bigl(\widehat P_q^{(1)}(\lambda,u),\,P_q^{(2)}(u)\bigr)
=\prod_{j=1}^{d_2}\widehat P_q^{(1)}(\lambda,\beta_j)
=\prod_{i=1}^{d_1}\prod_{j=1}^{d_2}(\lambda-\alpha_i\beta_j).
\]
故 \(R_q(\lambda)=\chi_{C_1\otimes C_2}(\lambda)\)。
由 Cayley--Hamilton 定理，
\(
R_q(C_1\otimes C_2)=0
\)，从而 \(R_q(E)S_q^{(\otimes)}\equiv 0\)。
由于 \(P_q^{(\otimes)}\) 按定义为该序列的首一最小递推特征多项式，必有
\(
P_q^{(\otimes)}\mid R_q
\)
于 \(\QQ[\lambda]\)。
证毕。
\end{proof}

\begin{theorem}[两因子张量化递推的精确结果式闭包]\label{thm:pom-twofactor-tensor-recurrence-exact-resultant-closure}
\leanverified{paper\_pom\_twofactor\_tensor\_recurrence\_exact\_resultant\_closure}
设
\[
S^{(1)}(m)=\sum_{a=1}^{d_1} c_{1,a}\,\alpha_{1,a}^{m},
\qquad
S^{(2)}(m)=\sum_{b=1}^{d_2} c_{2,b}\,\alpha_{2,b}^{m},
\]
其中 \(c_{1,a},c_{2,b}\neq 0\)，并定义逐项乘积序列
\[
S^{(1\boxtimes 2)}(m):=S^{(1)}(m)\,S^{(2)}(m).
\]
假设：
\begin{enumerate}
  \item \(\alpha_{1,1},\dots,\alpha_{1,d_1}\) 与 \(\alpha_{2,1},\dots,\alpha_{2,d_2}\) 全部非零，且在各自内部两两不同；
  \item 乘积谱无碰撞：
  \[
  \alpha_{1,a}\alpha_{2,b}=\alpha_{1,a'}\alpha_{2,b'}
  \Longrightarrow
  (a,b)=(a',b').
  \]
\end{enumerate}
令
\[
P_1(\lambda):=\prod_{a=1}^{d_1}(\lambda-\alpha_{1,a}),
\qquad
P_2(\lambda):=\prod_{b=1}^{d_2}(\lambda-\alpha_{2,b}),
\]
\[
\widetilde P_1(\lambda,u):=u^{d_1}P_1(\lambda/u),
\qquad
R_{1,2}(\lambda):=\operatorname{Res}_{u}\!\bigl(\widetilde P_1(\lambda,u),P_2(u)\bigr).
\]
则 \(S^{(1\boxtimes 2)}\) 的最小递推特征多项式恰为
\[
P_{1\boxtimes 2}(\lambda)=R_{1,2}(\lambda)
=
\prod_{a=1}^{d_1}\prod_{b=1}^{d_2}\bigl(\lambda-\alpha_{1,a}\alpha_{2,b}\bigr).
\]
特别地，
\[
\deg P_{1\boxtimes 2}=d_1d_2,
\qquad
\operatorname{rank}_H\!\bigl(S^{(1\boxtimes 2)}\bigr)=d_1d_2.
\]
\end{theorem}
\begin{proof}
由定义，
\[
S^{(1\boxtimes 2)}(m)
=
\sum_{a=1}^{d_1}\sum_{b=1}^{d_2}
\bigl(c_{1,a}c_{2,b}\bigr)\,
\bigl(\alpha_{1,a}\alpha_{2,b}\bigr)^m.
\]
由于 \(c_{1,a}c_{2,b}\neq 0\)，且乘积谱无碰撞，出现的指数模态
\[
\alpha_{1,a}\alpha_{2,b}
\qquad
(1\le a\le d_1,\ 1\le b\le d_2)
\]
两两不同。于是该序列的最小递推特征多项式必为
\[
\prod_{a=1}^{d_1}\prod_{b=1}^{d_2}
\bigl(\lambda-\alpha_{1,a}\alpha_{2,b}\bigr),
\]
故其次数与 Hankel 秩都等于 \(d_1d_2\)。

另一方面，
\[
\widetilde P_1(\lambda,u)
=
u^{d_1}P_1(\lambda/u)
=
\prod_{a=1}^{d_1}(\lambda-\alpha_{1,a}u).
\]
由于 \(P_2(u)=\prod_{b=1}^{d_2}(u-\alpha_{2,b})\) 首一，结果式的根值公式给出
\[
R_{1,2}(\lambda)
=
\prod_{b=1}^{d_2}\widetilde P_1(\lambda,\alpha_{2,b})
=
\prod_{a=1}^{d_1}\prod_{b=1}^{d_2}
\bigl(\lambda-\alpha_{1,a}\alpha_{2,b}\bigr).
\]
故 \(R_{1,2}\) 与上述最小递推特征多项式一致。证毕。
\end{proof}

\begin{theorem}[多因子张量化递推的严格乘法化]\label{thm:pom-multifactor-tensor-recurrence-strict-multiplicativity}
设 \(r\ge 2\)，并有 \(r\) 个序列
\[
S^{(i)}(m)=\sum_{a_i=1}^{d_i} c_{i,a_i}\,\alpha_{i,a_i}^{m}
\qquad
(1\le i\le r),
\]
其中 \(c_{i,a_i}\neq 0\)。假设：
\begin{enumerate}
  \item 对每个 \(i\)，多项式
  \[
  P_i(\lambda):=\prod_{a_i=1}^{d_i}(\lambda-\alpha_{i,a_i})
  \]
  的根全为非零单根；
  \item 全乘积无碰撞：
  \[
  \prod_{i=1}^{r}\alpha_{i,a_i}
  =
  \prod_{i=1}^{r}\alpha_{i,b_i}
  \Longrightarrow
  a_i=b_i\quad\text{对一切 }i.
  \]
\end{enumerate}
定义 \(r\)-重逐项乘积序列
\[
S^{(\boxtimes r)}(m):=\prod_{i=1}^{r}S^{(i)}(m).
\]
递归定义迭代结果式
\[
R^{[1]}(\lambda):=P_1(\lambda),
\]
\[
R^{[k]}(\lambda):=
\operatorname{Res}_{u}\!\Bigl(
u^{\deg R^{[k-1]}}R^{[k-1]}(\lambda/u),\,P_k(u)
\Bigr)
\qquad (2\le k\le r).
\]
则 \(S^{(\boxtimes r)}\) 的最小递推特征多项式满足
\[
P_{\boxtimes r}(\lambda)=R^{[r]}(\lambda)
=
\prod_{a_1=1}^{d_1}\cdots\prod_{a_r=1}^{d_r}
\Bigl(\lambda-\prod_{i=1}^{r}\alpha_{i,a_i}\Bigr).
\]
从而
\[
\deg P_{\boxtimes r}=\prod_{i=1}^{r}d_i,
\qquad
\operatorname{rank}_H\!\bigl(S^{(\boxtimes r)}\bigr)=\prod_{i=1}^{r}d_i.
\]
\leanverified{paper\_pom\_multifactor\_tensor\_recurrence\_strict\_multiplicativity}
\end{theorem}
\begin{proof}
对 \(r=2\) 时，结论即为定理
\ref{thm:pom-twofactor-tensor-recurrence-exact-resultant-closure}。
现对 \(r\) 作归纳。设结论对 \(r-1\) 已成立，则
\[
S^{(\boxtimes (r-1))}(m)
=
\sum_{\mathbf a} C_{\mathbf a}\,\beta_{\mathbf a}^{m},
\qquad
\beta_{\mathbf a}:=\prod_{i=1}^{r-1}\alpha_{i,a_i},
\]
其中 \(\mathbf a=(a_1,\dots,a_{r-1})\)，且所有 \(\beta_{\mathbf a}\) 两两不同，并且
\[
P_{\boxtimes (r-1)}(\lambda)
=
\prod_{\mathbf a}(\lambda-\beta_{\mathbf a})
=
R^{[r-1]}(\lambda).
\]
再与第 \(r\) 个序列逐项相乘，得到
\[
S^{(\boxtimes r)}(m)
=
\sum_{\mathbf a}\sum_{a_r=1}^{d_r}
\bigl(C_{\mathbf a}c_{r,a_r}\bigr)\,
\bigl(\beta_{\mathbf a}\alpha_{r,a_r}\bigr)^m.
\]
由全乘积无碰撞假设，所有
\[
\beta_{\mathbf a}\alpha_{r,a_r}
=
\prod_{i=1}^{r}\alpha_{i,a_i}
\]
两两不同，因此 \(S^{(\boxtimes r)}\) 的最小递推特征多项式恰为
\[
\prod_{\mathbf a}\prod_{a_r=1}^{d_r}
\bigl(\lambda-\beta_{\mathbf a}\alpha_{r,a_r}\bigr).
\]
另一方面，把定理
\ref{thm:pom-twofactor-tensor-recurrence-exact-resultant-closure}
应用于序列 \(S^{(\boxtimes(r-1))}\) 与 \(S^{(r)}\)，便得
\[
R^{[r]}(\lambda)
=
\prod_{\mathbf a}\prod_{a_r=1}^{d_r}
\bigl(\lambda-\beta_{\mathbf a}\alpha_{r,a_r}\bigr).
\]
故 \(P_{\boxtimes r}=R^{[r]}\)。次数公式与 Hankel 秩公式随之立即得到。证毕。
\end{proof}
\begin{corollary}[非碰撞制度下的 Hankel 复杂度严格乘法性]\label{cor:pom-tensor-fold-hankel-rank-multiplicativity}
\leanverified{paper\_pom\_tensor\_fold\_hankel\_rank\_multiplicativity}
沿用定理 \ref{thm:pom-tensor-fold-resultant-recursion} 的设定，并记
\(
\operatorname{rank}_H(\cdot)
\)
为序列的 Hankel 秩（等价地，其最小递推阶数）。
若进一步满足：
\begin{enumerate}
  \item \(P_q^{(1)}\) 与 \(P_q^{(2)}\) 在 \(\overline{\QQ}\) 上均无重根；
  \item 按其根展开的谱系数均非零，即存在
  \[
  S_q^{(1)}(m)=\sum_{i=1}^{d_1}a_i\alpha_i^m,
  \qquad
  S_q^{(2)}(m)=\sum_{j=1}^{d_2}b_j\beta_j^m
  \qquad(a_i,b_j\neq 0) ;
  \]
  \item 谱乘积无碰撞：
  \[
  \alpha_i\beta_j=\alpha_{i'}\beta_{j'}
  \Longrightarrow
  (i,j)=(i',j').
  \]
\end{enumerate}
则
\[
\deg P_q^{(\otimes)}=d_1d_2,
\qquad
\operatorname{rank}_H\!\bigl(S_q^{(\otimes)}\bigr)
=\operatorname{rank}_H\!\bigl(S_q^{(1)}\bigr)\,
\operatorname{rank}_H\!\bigl(S_q^{(2)}\bigr).
\]
\end{corollary}
\begin{proof}
由简单根分解与系数非零性，
\[
S_q^{(\otimes)}(m)
=S_q^{(1)}(m)S_q^{(2)}(m)
=\sum_{i=1}^{d_1}\sum_{j=1}^{d_2}(a_ib_j)\,(\alpha_i\beta_j)^m.
\]
由谱乘积无碰撞，\(\{\alpha_i\beta_j\}_{i,j}\) 两两不同。
对特征零域上的两两不同标量 \(\gamma_1,\dots,\gamma_N\)，指数序列 \(\gamma_k^m\)（\(0\le m\le N-1\)）线性无关；否则 Vandermonde 行列式
\[
\det(\gamma_k^{\,m})_{0\le m\le N-1,\ 1\le k\le N}
=\prod_{1\le k<\ell\le N}(\gamma_\ell-\gamma_k)
\]
将为零，矛盾。
因此 \(\bigl(S_q^{(\otimes)}(m)\bigr)\) 的最小递推阶数至少为 \(d_1d_2\)。

另一方面，定理 \ref{thm:pom-tensor-fold-resultant-recursion} 表明
\(
P_q^{(\otimes)}\mid R_q
\)，而 \(\deg R_q=d_1d_2\)。
故 \(\deg P_q^{(\otimes)}\le d_1d_2\)。
两侧合并即得 \(\deg P_q^{(\otimes)}=d_1d_2\)。
再由 Hankel 秩与最小递推阶数的一致性，即得所示严格乘法律。证毕。
\end{proof}

\begin{theorem}[张量化矩谱的 Hadamard--\texorpdfstring{\(\zeta\)}{zeta} 精确极点律与压力可加性]\label{thm:pom-tensor-hadamard-pole-law-pressure-additivity}
在定理 \ref{thm:pom-multifactor-tensor-recurrence-strict-multiplicativity} 的设定下，对每个 \(1\le i\le r\) 定义普通生成函数
\[
F_i(z):=\sum_{m\ge 0}S^{(i)}(m)z^m.
\]
则其 \(r\)-重 Hadamard 乘积
\[
F_{\boxtimes r}(z):=(F_1\odot\cdots\odot F_r)(z)
=
\sum_{m\ge 0}\Bigl(\prod_{i=1}^{r}S^{(i)}(m)\Bigr)z^m
\]
满足精确分解
\[
F_{\boxtimes r}(z)
=
\sum_{a_1,\dots,a_r}
\frac{\prod_{i=1}^{r}c_{i,a_i}}
{1-\bigl(\prod_{i=1}^{r}\alpha_{i,a_i}\bigr)z}.
\]
因此：
\begin{enumerate}
  \item \(F_{\boxtimes r}\) 是有理函数，其全部极点均为单极点；
  \item 极点集合精确为
  \[
  \left\{
  \Bigl(\prod_{i=1}^{r}\alpha_{i,a_i}\Bigr)^{-1}
  \right\}_{a_1,\dots,a_r};
  \]
  \item 若记各因子的谱半径与对应压力为
  \[
  \rho_i:=\max_{a_i}|\alpha_{i,a_i}|,
  \qquad
  \mathcal P_i:=\log \rho_i,
  \]
  则
  \[
  \rho_{\boxtimes r}
  :=
  \max_{a_1,\dots,a_r}
  \left|
  \prod_{i=1}^{r}\alpha_{i,a_i}
  \right|
  =
  \prod_{i=1}^{r}\rho_i,
  \]
  \[
  \mathcal P_{\boxtimes r}
  :=
  \log \rho_{\boxtimes r}
  =
  \sum_{i=1}^{r}\mathcal P_i.
  \]
\end{enumerate}
\end{theorem}
\begin{proof}
由定义，
\[
\prod_{i=1}^{r}S^{(i)}(m)
=
\sum_{a_1,\dots,a_r}
\Bigl(\prod_{i=1}^{r}c_{i,a_i}\Bigr)
\Bigl(\prod_{i=1}^{r}\alpha_{i,a_i}\Bigr)^m.
\]
对 \(m\ge 0\) 求和即得
\[
F_{\boxtimes r}(z)
=
\sum_{a_1,\dots,a_r}\sum_{m\ge 0}
\Bigl(\prod_{i=1}^{r}c_{i,a_i}\Bigr)
\Bigl(\bigl(\prod_{i=1}^{r}\alpha_{i,a_i}\bigr)z\Bigr)^m
=
\sum_{a_1,\dots,a_r}
\frac{\prod_{i=1}^{r}c_{i,a_i}}
{1-\bigl(\prod_{i=1}^{r}\alpha_{i,a_i}\bigr)z}.
\]
这就给出了所述部分分式展开。

由定理 \ref{thm:pom-multifactor-tensor-recurrence-strict-multiplicativity} 的全乘积无碰撞假设，所有分母根彼此不同，故全部极点均为单极点，且极点集合正是所示乘积倒数集合。

最后，
\[
\max_{a_1,\dots,a_r}
\left|
\prod_{i=1}^{r}\alpha_{i,a_i}
\right|
=
\prod_{i=1}^{r}\max_{a_i}|\alpha_{i,a_i}|
=
\prod_{i=1}^{r}\rho_i.
\]
取对数即得压力可加性
\[
\mathcal P_{\boxtimes r}
=
\log \rho_{\boxtimes r}
=
\sum_{i=1}^{r}\log \rho_i
=
\sum_{i=1}^{r}\mathcal P_i.
\]
证毕。
\end{proof}

\begin{theorem}[张量化谱根的 Galois 刚性：次数恰为乘积]\label{thm:pom-tensor-product-root-galois-degree-rigidity}
\leanverified{paper\_pom\_tensor\_product\_root\_galois\_degree\_rigidity}
在定理 \ref{thm:pom-multifactor-tensor-recurrence-strict-multiplicativity} 的设定下，进一步假设：
\begin{enumerate}
  \item 每个 \(P_i\in\QQ[\lambda]\) 在 \(\QQ\) 上不可约，次数为 \(d_i\)；
  \item 若记 \(L_i/\QQ\) 为 \(P_i\) 的分裂域、\(G_i:=\Gal(L_i/\QQ)\)，则 \(\{L_i\}_{i=1}^{r}\) 在 \(\QQ\) 上相互线性无交；等价地，自然限制映射
  \[
  \Gal(L_1\cdots L_r/\QQ)\longrightarrow \prod_{i=1}^{r}G_i
  \]
  为同构；
  \item 每个 \(G_i\) 在 \(P_i\) 的根集上传递作用；
  \item 全乘积无碰撞成立。
\end{enumerate}
取任意固定根 \(\alpha_{i,1}\in L_i\)，并设
\[
\gamma:=\prod_{i=1}^{r}\alpha_{i,1}.
\]
则
\[
[\QQ(\gamma):\QQ]=\prod_{i=1}^{r}d_i.
\]
并且 \(\gamma\) 的最小多项式恰等于定理
\ref{thm:pom-multifactor-tensor-recurrence-strict-multiplicativity}
中的迭代结果式 \(R^{[r]}(\lambda)\)。
\end{theorem}
\begin{proof}
记 \(L:=L_1\cdots L_r\)。由假设，限制映射给出同构
\[
\Gal(L/\QQ)\cong \prod_{i=1}^{r}G_i.
\]
每个 \(G_i\) 在 \(P_i\) 的根集上传递，故直积群在根元组集合
\[
\mathcal O
:=
\{(\alpha_{1,a_1},\dots,\alpha_{r,a_r})\}
\]
上传递，其大小为
\[
|\mathcal O|=\prod_{i=1}^{r}d_i.
\]
考虑乘法映射
\[
\mu:\mathcal O\to \overline{\QQ},
\qquad
(\alpha_{1,a_1},\dots,\alpha_{r,a_r})
\longmapsto
\prod_{i=1}^{r}\alpha_{i,a_i}.
\]
由全乘积无碰撞，\(\mu\) 在 \(\mathcal O\) 上单射，因此 \(\gamma\) 在 \(\Gal(L/\QQ)\) 下的共轭轨道大小恰为
\[
\prod_{i=1}^{r}d_i.
\]
由于 \(L/\QQ\) 可分，\(\gamma\) 的最小多项式次数等于其共轭轨道大小，故
\[
[\QQ(\gamma):\QQ]=\prod_{i=1}^{r}d_i.
\]

另一方面，定理
\ref{thm:pom-multifactor-tensor-recurrence-strict-multiplicativity}
已表明迭代结果式 \(R^{[r]}(\lambda)\) 的根恰为全部乘积
\[
\prod_{i=1}^{r}\alpha_{i,a_i},
\]
其次数也是 \(\prod_{i=1}^{r}d_i\)。由于 \(\gamma\) 是其一根，而 \(\gamma\) 的最小多项式次数已经达到该值，故 \(\gamma\) 的最小多项式只能等于 \(R^{[r]}(\lambda)\)。证毕。
\end{proof}

\endinput


\begin{theorem}[\texorpdfstring{$q$}{q}-重同步直积的置换对称与碰撞核的内禀对易子]\label{thm:pom-momq-permutation-symmetry}
沿用定理 \ref{thm:pom-fiber-multiplicity-matrix-product} 与命题 \ref{prop:pom-momq-tensor-kernel-construction} 的记号。
令 \(V:=\CC^{\mathcal S}\)，并令 \(T_b:V\to V\)（\(b\in\{0,1\}\)）为相应的线性算子；对整数 \(q\ge 2\) 令 \(W:=V^{\otimes q}\)。
定义置换表示
\[
\pi:S_q\to \mathrm{GL}(W),
\qquad
\pi(\sigma)(v_1\otimes\cdots\otimes v_q)
=v_{\sigma^{-1}(1)}\otimes\cdots\otimes v_{\sigma^{-1}(q)}.
\]
则对每个 \(b\in\{0,1\}\) 与 \(\sigma\in S_q\) 成立交换关系
\[
\pi(\sigma)\,T_b^{\otimes q}=T_b^{\otimes q}\,\pi(\sigma).
\]
从而
\[
I_2\otimes \pi(S_q)\ \subseteq\ \mathrm{Comm}\bigl(A_q^{\mathrm{ten}}\bigr),
\]
即置换对称作为 \(q\)-重同步直积的结构性对易子，内禀地进入张量分解型碰撞核 \(A_q^{\mathrm{ten}}=G_0\otimes T_0^{\otimes q}+G_1\otimes T_1^{\otimes q}\)。
\end{theorem}
\begin{proof}
对任意 \(v_1,\dots,v_q\in V\)，由张量积函子性
\[
T_b^{\otimes q}(v_1\otimes\cdots\otimes v_q)=T_bv_1\otimes\cdots\otimes T_bv_q
\]
可见：先对张量因子作置换、再施加 \(T_b^{\otimes q}\) 与先施加 \(T_b^{\otimes q}\)、再置换完全一致，故 \(\pi(\sigma)\) 与 \(T_b^{\otimes q}\) 对易。
于是 \(I_2\otimes \pi(\sigma)\) 分别与 \(G_0\otimes T_0^{\otimes q}\) 与 \(G_1\otimes T_1^{\otimes q}\) 对易，从而与其和 \(A_q^{\mathrm{ten}}\) 对易。
\end{proof}
\leanverified{paper\_pom\_momq\_permutation\_symmetry}

\begin{theorem}[Schur--Weyl 最大性：不可约转移代数下的置换对易子刻画]\label{thm:pom-momq-schur-weyl-commutant}
沿用上述记号，并令
\[
\mathcal A:=\mathrm{Alg}_{\CC}\langle T_0,T_1\rangle\ \subseteq\ \mathrm{End}(V)
\]
为由逆码转移算子生成的复代数。
若 \(\mathcal A\) 在 \(V\) 上作用不可约，则对每个整数 \(q\ge 2\) 有
\[
\mathrm{Comm}_{\mathrm{End}(V^{\otimes q})}\!\bigl(\mathcal A^{\otimes q}\bigr)
\;=\;
\pi\bigl(\CC[S_q]\bigr),
\]
亦即：在张量因子上与 \(T_0^{\otimes q},T_1^{\otimes q}\)（等价地与 \(\mathcal A^{\otimes q}\)）同时对易的全部算子恰由置换作用生成。
\end{theorem}
\begin{proof}
由 Burnside 定理，\(\mathcal A\) 在 \(V\) 上不可约推出 \(\mathcal A=\mathrm{End}(V)\)。
于是 \(\mathcal A^{\otimes q}=\mathrm{End}(V)^{\otimes q}\) 在 \(V^{\otimes q}\) 上的对易代数由 Schur--Weyl 对偶性给出，正是置换表示的像 \(\pi(\CC[S_q])\)。
\end{proof}
\leanverified{paper\_pom\_momq\_schur\_weyl\_commutant}


\paragraph{不交张量底座：\texorpdfstring{$M^{\otimes q}$}{M^⊗q} 的谱显式与四阶整数临界}
令
\(
M:=G_0+G_1=\bigl(\begin{smallmatrix}1&1\\[2pt]1&0\end{smallmatrix}\bigr)
\)。
该两态矩阵不仅编码黄金均值语法骨架，亦给出一个更原初的张量硬核对象：不交图的邻接算子。

\begin{theorem}[两态骨架的中心化半环与 Fibonacci 融合律]\label{thm:pom-centralizer-semiring-fibonacci-fusion}
令
\[
M:=\begin{pmatrix}1&1\\[2pt]1&0\end{pmatrix},
\qquad
\mathcal C_+(M):=\{A\in M_2(\ZZ_{\ge 0}):\ AM=MA\}.
\]
则存在半环同构
\[
\boxed{\;
\mathcal C_+(M)\ \cong\ \ZZ_{\ge 0}[\tau]/(\tau^2-\tau-1),\;}
\]
其同构由 \(\tau\mapsto M\) 实现。
特别地，\(\mathcal C_+(M)\) 由 \(I\) 与 \(M\) 生成并满足
\[
M^2=M+I,
\]
从而其乘法结构与 Fibonacci 融合环的去范畴化融合律同型。
\leanverified{paper\_pom\_centralizer\_semiring\_fibonacci\_fusion}
\end{theorem}
\begin{proof}
任取 \(A=\bigl(\begin{smallmatrix}p&q\\ r&s\end{smallmatrix}\bigr)\in M_2(\ZZ)\)。直接计算
\[
AM=\begin{pmatrix}p+q&p\\ r+s&r\end{pmatrix},
\qquad
MA=\begin{pmatrix}p+r&q+s\\ p&q\end{pmatrix}.
\]
令 \(AM=MA\) 得方程组
\(
q=r,\ p=q+s
\)。
因此
\[
A=\begin{pmatrix}q+s&q\\ q&s\end{pmatrix}=s\,I+q\,M.
\]
若进一步要求 \(A\in M_2(\ZZ_{\ge 0})\)，则 \(q,s\ge 0\)，于是
\(
\mathcal C_+(M)=\{\,aI+bM:\ a,b\in\ZZ_{\ge 0}\,\}
\)。
另一方面，由直接乘法得 \(M^2=M+I\)。因此映射 \(\ZZ_{\ge 0}[\tau]\to\mathcal C_+(M)\)（\(\tau\mapsto M\)）为满射半环同态，其核包含 \((\tau^2-\tau-1)\)，并由表示 \(aI+bM\) 的唯一性知核恰为该主理想，故得所示同构。证毕。
\end{proof}

\begin{corollary}[Frobenius--Perron 维数的黄金参数]\label{cor:pom-centralizer-fpdim-golden}
令 \(\mathrm{FPdim}(\tau)\) 以谱半径刻画，则
\[
\mathrm{FPdim}(\tau)=\rho(M)=\varphi,
\]
其中 \(\varphi=(1+\sqrt5)/2\)。
\end{corollary}
\begin{proof}
\(M\) 的特征多项式为 \(\lambda^2-\lambda-1\)，其 Perron 根为 \(\varphi\)。证毕。
\leanverified{paper\_centralizer\_fpdim\_golden}
\end{proof}

\begin{proposition}[中央化子中的二次范数与特征 \texorpdfstring{$5$}{5} 的平方像现象]\label{prop:pom-centralizer-det-norm-mod5}
令
\[
M=\begin{pmatrix}1&1\\[2pt]1&0\end{pmatrix},
\qquad
C(M):=\{A\in M_2(\ZZ):\ AM=MA\}.
\]
则 \(C(M)=\ZZ[M]\)，且对任意 \(A\in \ZZ[M]\) 唯一可写为 \(A=aI+bM\)（\(a,b\in\ZZ\)），并满足
\begin{equation}\label{eq:pom-centralizer-det-norm-form}
\det(A)=a^2+ab-b^2.
\end{equation}
此外，在模 \(5\) 意义下恒有
\begin{equation}\label{eq:pom-centralizer-det-square-mod5}
\det(A)\equiv (a+3b)^2\pmod 5,
\end{equation}
从而 \(\det(\ZZ[M]\otimes \FF_5)\) 的像包含于 \(\FF_5\) 的平方元集合。
反之，对任意素数 \(p\neq 5\)，\(\det(\ZZ[M]\otimes \FF_p)\) 的像为整个 \(\FF_p\)。
\end{proposition}
\begin{proof}
由定理 \ref{thm:pom-centralizer-semiring-fibonacci-fusion} 的整数版本，\(C(M)=\ZZ[M]\) 且任意 \(A\in\ZZ[M]\) 唯一写作 \(A=aI+bM\)。
直接计算
\[
A=aI+bM
=\begin{pmatrix}a+b&b\\[2pt]b&a\end{pmatrix},
\qquad
\det(A)=(a+b)a-b^2=a^2+ab-b^2,
\]
得到 \eqref{eq:pom-centralizer-det-norm-form}。
\par
在 \(\FF_5\) 中有 \(-1\equiv 4\)，故
\(
a^2+ab-b^2\equiv a^2+ab+4b^2
\)。
取 \(k\equiv 3\pmod 5\)，则 \(2k\equiv 1\pmod 5\) 且 \(k^2\equiv 9\equiv 4\pmod 5\)，于是
\[
(a+kb)^2\equiv a^2+2kab+k^2b^2\equiv a^2+ab+4b^2\pmod 5,
\]
即得 \eqref{eq:pom-centralizer-det-square-mod5}。
\par
对 \(p\neq 5\)，二次多项式 \(\tau^2-\tau-1\) 的判别式为 \(5\)，在 \(\FF_p\) 上不为零，故
\(
\FF_p[\tau]/(\tau^2-\tau-1)
\)
为二次 \'{e}tale \(\FF_p\)-代数（分裂为 \(\FF_p\times\FF_p\) 或为二次扩张 \(\FF_{p^2}\)）。
在同构 \(\ZZ[M]\otimes\FF_p\cong \FF_p[\tau]/(\tau^2-\tau-1)\) 下，\(\det(aI+bM)\) 恰为该二次 \'{e}tale 代数到 \(\FF_p\) 的范数。
而有限域扩张的范数映射在乘法群上满射，分裂情形亦为两分量乘积之满射；故 \(\det\) 的像为 \(\FF_p\) 的全体元素。
\leanverified{centralizer\_det\_formula}
\leanverified{det\_mod5\_is\_square}
\leanverified{golden\_ratio\_discriminant}
\leanverified{golden\_poly\_mod5\_double\_root}
\leanverified{paper\_pom\_centralizer\_det\_mod5}
\leanverified{paper\_pom\_centralizer\_det\_norm\_mod5}
\end{proof}

\begin{corollary}[判别式锁定的唯一坏素数：\texorpdfstring{$p=5$}{p=5} 的骨架退化]\label{cor:pom-centralizer-unique-bad-prime-5}
在中央化子骨架 \(\ZZ[M]\cong \ZZ[\tau]/(\tau^2-\tau-1)\) 上，模素约化 \(\ZZ[M]\otimes\FF_p\) 为非 \'{e}tale（非约化）当且仅当 \(p=5\)；
等价地，特征 \(5\) 是该二次骨架在纯判别式层面发生退化的唯一算术奇点。
\end{corollary}
\begin{proof}
\(\tau^2-\tau-1\) 的判别式为 \(5\)。
对 \(p\neq 5\)，\(\FF_p[\tau]/(\tau^2-\tau-1)\) 为 \'{e}tale 代数；
对 \(p=5\)，有 \(\tau^2-\tau-1\equiv(\tau-3)^2\) 于 \(\FF_5\)，从而约化环非约化并带幂零元。
由 \(\ZZ[M]\cong \ZZ[\tau]/(\tau^2-\tau-1)\) 得证。
\end{proof}
\leanverified{paper\_pom\_centralizer\_unique\_bad\_prime\_5}


\begin{proposition}[不交图的张量幂刻画与 \texorpdfstring{$\QQ(\sqrt5)$}{Q(sqrt5)} 的普适谱底座]\label{prop:pom-disjointness-graph-tensor-power-spectrum}
令 \(\mathcal{D}_q\) 为顶点集 \(2^{[q]}\) 上的不交图：对 \(S,T\subseteq[q]\)，令 \(S\sim T\) 当且仅当 \(S\cap T=\varnothing\)。
在自然基（以 \(\{0,1\}^q\cong 2^{[q]}\) 标号）下，\(\mathcal{D}_q\) 的邻接算子严格等于
\[
\boxed{\;A(\mathcal{D}_q)=M^{\otimes q}\;}
\qquad\Bigl(M=\begin{pmatrix}1&1\\[2pt]1&0\end{pmatrix}\Bigr).
\]
从而其谱完全显式：若 \(\varphi=(1+\sqrt5)/2\)，则全部特征值为
\[
(-1)^k\varphi^{\,q-2k}\qquad(0\le k\le q),
\]
且代数重数为 \(\binom{q}{k}\)。
因此不交张量底座的谱分解在每个 \(q\) 上均由 \(\QQ(\sqrt5)\) 统摄。
\end{proposition}
\begin{proof}
将 \(S\subseteq[q]\) 识别为指示向量 \(s\in\{0,1\}^q\)，则 \(S\cap T=\varnothing\) 等价于对每个坐标 \(i\)，不允许出现局部对 \((s_i,t_i)=(1,1)\)。
故邻接矩阵按坐标分解为 \(q\) 个独立的 \(2\times 2\) 局部允许矩阵之 Kronecker 乘积，即 \(A(\mathcal{D}_q)=M^{\otimes q}\)。
又 \(M\) 的特征值为 \(\varphi\) 与 \(-\varphi^{-1}\)，故张量幂的特征值为 \(\varphi^{q-k}(-\varphi^{-1})^k=(-1)^k\varphi^{q-2k}\)，重数由二项式展开给出。证毕。
\end{proof}
\leanverified{paper\_pom\_disjointness\_graph\_tensor\_power\_spectrum}

\begin{corollary}[\texorpdfstring{$q=4$}{q=4} 为不交图层面 Ramanujan 比率的唯一整数临界点]\label{cor:pom-disjointness-ramanujan-q4-critical}
令 \(r_q^{\mathrm{disj}}:=\rho(A(\mathcal{D}_q))\) 为 Perron 根，并令 \(\Lambda_q^{\mathrm{disj}}\) 为非 Perron 最大模特征值。
则
\[
r_q^{\mathrm{disj}}=\varphi^{\,q},\qquad \Lambda_q^{\mathrm{disj}}=\varphi^{\,q-2},\qquad
\frac{\Lambda_q^{\mathrm{disj}}}{\sqrt{r_q^{\mathrm{disj}}}}=\varphi^{\,q/2-2}.
\]
因此对整数 \(q\le 3\) 有 \(\Lambda_q^{\mathrm{disj}}<\sqrt{r_q^{\mathrm{disj}}}\)，对 \(q=4\) 有临界等式 \(\Lambda_4^{\mathrm{disj}}=\sqrt{r_4^{\mathrm{disj}}}\)，而对整数 \(q\ge 5\) 必有 \(\Lambda_q^{\mathrm{disj}}>\sqrt{r_q^{\mathrm{disj}}}\)。
\end{corollary}
\begin{proof}
由命题 \ref{prop:pom-disjointness-graph-tensor-power-spectrum}，\(r_q^{\mathrm{disj}}=\varphi^q\)，且非 Perron 模最大者来自 \(k=1\) 分支，故 \(\Lambda_q^{\mathrm{disj}}=\varphi^{q-2}\)。
代入即得比率闭式，并由 \(\varphi>1\) 推出临界点的唯一性。证毕。
\end{proof}
\leanverified{paper\_pom\_disjointness\_ramanujan\_q4\_critical}

\begin{conclusion}[Kronecker 不交核的量子化临界：\texorpdfstring{$S_q$}{Sq}-等变压缩的临界 \(q\) 被最小权阶量子化为 \(4k_0\)]\label{con:pom-disjointness-sq-equivariant-critical-quantization}
沿命题 \ref{prop:pom-disjointness-graph-tensor-power-spectrum}，令
\(
M=\bigl(\begin{smallmatrix}1&1\\[2pt]1&0\end{smallmatrix}\bigr)
\)
并考虑 \(M^{\otimes q}\) 在 \((\RR^2)^{\otimes q}\) 上的作用。
设 \(W\subset(\RR^2)^{\otimes q}\) 为任意 \(S_q\)-不变子空间，且包含 Perron 模态（等价地，包含 \(M^{\otimes q}\) 的 Perron 特征向量）。
令
\[
T:=\mathrm{Proj}_W(M^{\otimes q})\mathrm{Proj}_W
\]
为相应的 \(S_q\)-等变压缩算子，并令 \(k_0\) 为 \(W\) 所包含的最小非零“权阶”：即最小的 \(k\ge 1\) 使得 \(W\) 在特征值 \((-1)^k\varphi^{q-2k}\) 的特征空间上的投影非零。
则压缩算子仍满足 \(r=\rho(T)=\varphi^q\)，并且其最大非 Perron 模长精确为
\[
\Lambda=\varphi^{q-2k_0},
\qquad
\frac{\Lambda}{\sqrt r}=\varphi^{\,q/2-2k_0}.
\]
因此 Ramanujan 临界条件 \(\Lambda\le \sqrt r\) 等价于 \(q\le 4k_0\)，临界点被量子化为
\[
\boxed{\;q_c=4k_0\;}
\]
（在整数 \(q\) 上即表现为离散阈值序列）。
特别地，\(k_0=1\) 对应的最小非平凡压缩给出 \(q=4\) 的临界性。
\end{conclusion}

\begin{proposition}[\texorpdfstring{$S_q$}{Sq}-对称压缩的 \texorpdfstring{$(q+1)$}{q+1} 维闭式矩阵与四阶整数 \(7\) 的谱不变量]\label{prop:pom-disjointness-symmetric-compression-bq}
\leanverified{paper\_pom\_disjointness\_symmetric\_compression\_bq}
令 \(B_q\) 为 \(\mathcal{D}_q\) 在 \(S_q\) 轨道空间（按子集大小分层）上的诱导算子。
则 \(B_q\) 为 \((q+1)\times(q+1)\) 矩阵；在以“大小层”标号的基下，其系数满足显式闭式
\[
\boxed{\;B_q(i,j)=\binom{q-j}{i}\;}\qquad(0\le i,j\le q).
\]
特别地，当 \(q=4\) 时其特征多项式严格因式分解为
\[
(\lambda-1)(\lambda^2-7\lambda+1)(\lambda^2+3\lambda+1),
\]
其五个特征值分别为
\(
1
\)、\(
\frac{7\pm 3\sqrt5}{2}=\varphi^{\pm 4}
\)、\(
\frac{-3\pm\sqrt5}{2}=-\varphi^{\pm 2}
\)。
因此四阶层面的“系数 \(7\)”在不交张量底座的最小非平凡对称压缩中已作为谱不变量出现。
\end{proposition}
\begin{proof}
对固定 \(T\subseteq[q]\) 且 \(|T|=j\)，与之不交的 \(i\)-元子集数恰为从补集 \([q]\setminus T\)（大小 \(q-j\)）中取 \(i\) 元子集的计数 \(\binom{q-j}{i}\)，故 \(B_q(i,j)=\binom{q-j}{i}\)。
当 \(q=4\) 时对该 \(5\times 5\) 矩阵直接计算特征多项式并因式分解即可得到所示闭式。证毕。
\end{proof}

% (User input window)

\begin{proposition}[二项权内积下的严格对称化：\texorpdfstring{$B_q$}{Bq} 的带权自伴随性]\label{prop:pom-bq-weighted-selfadjoint-binomial}
\leanverified{paper\_pom\_bq\_weighted\_selfadjoint\_binomial}
沿命题 \ref{prop:pom-disjointness-symmetric-compression-bq} 的记号，令 \(B_q\in M_{q+1}(\ZZ)\) 为满足 \(B_q(i,j)=\binom{q-j}{i}\) 的矩阵。
定义对角权矩阵
\[
W_q:=\mathrm{diag}\Bigl(\binom{q}{0}^{-1},\binom{q}{1}^{-1},\dots,\binom{q}{q}^{-1}\Bigr).
\]
则成立严格恒等式
\[
\boxed{\;W_qB_q=B_q^{\mathsf T}W_q\;}
\]
从而 \(B_q\) 在加权内积 \(\langle u,v\rangle_{W_q}:=u^{\mathsf T}W_qv\) 下自伴随。
等价地，令
\[
D_q:=\mathrm{diag}\Bigl(\sqrt{\binom{q}{0}},\dots,\sqrt{\binom{q}{q}}\Bigr),
\qquad
J_q:=D_q^{-1}B_qD_q,
\]
则 \(J_q\) 为实对称矩阵。
因此 \(B_q\) 的全部谱为实数，并可取一组 \(W_q\)-正交特征向量基。
\end{proposition}
\begin{proof}
对任意 \(0\le i,j\le q\)，由 \(B_q(i,j)=\binom{q-j}{i}\) 得
\[
(W_qB_q)(i,j)=\binom{q}{i}^{-1}\binom{q-j}{i},
\qquad
(B_q^{\mathsf T}W_q)(i,j)=\binom{q-i}{j}\binom{q}{j}^{-1}.
\]
当 \(i+j>q\) 时两者皆为 \(0\)。
当 \(i+j\le q\) 时直接化简得
\[
\binom{q}{i}^{-1}\binom{q-j}{i}
=\frac{(q-i)!(q-j)!}{q!(q-i-j)!}
=\binom{q}{j}^{-1}\binom{q-i}{j},
\]
故 \(W_qB_q=B_q^{\mathsf T}W_q\)。
另一方面，由 \(W_q=D_q^{-2}\)，上式等价于 \(J_q^{\mathsf T}=J_q\)，从而 \(J_q\) 可被正交对角化；回到相似变换即可得到 \(B_q\) 的谱实性与 \(W_q\)-正交特征基的存在性。证毕。
\end{proof}

% (User input window)

\paragraph{Lucas 平方谱分解与闭式迹公式：\texorpdfstring{$B_q$}{Bq} 的完全可计算周期动力学}
令
\(
M=\bigl(\begin{smallmatrix}1&1\\[2pt]1&0\end{smallmatrix}\bigr)
\)
并记 Lucas 数列为 \(L_0=2,L_1=1\) 且 \(L_{n+1}=L_n+L_{n-1}\)。

\begin{proposition}[\texorpdfstring{$B_q=\mathrm{Sym}^q(M)$}{Bq=Sym^q(M)} 与谱闭式]\label{prop:pom-bq-is-symq-and-spectrum}
\leanverified{paper\_pom\_bq\_is\_symq\_and\_spectrum}
对任意整数 \(q\ge 1\)，对称压缩核 \(B_q\) 与 \(M\) 的 \(q\)-次对称幂同构：
\[
\boxed{\;B_q\ \cong\ \mathrm{Sym}^q(M)\in \mathrm{GL}_{q+1}(\ZZ)\;}
\]
从而其全部特征值严格为
\[
\boxed{\;\mu_k=(-1)^k\varphi^{\,q-2k}\qquad(0\le k\le q)\;}
\]
且每个模态在 \(B_q\) 上的代数重数均为 \(1\)。
\end{proposition}
\begin{proof}
令 \(V=\RR^2\) 并取基向量 \(e_0,e_1\)。
用标准同构将 \(\mathrm{Sym}^q(V)\) 表示为齐次次数为 \(q\) 的二元多项式空间：形式变量 \(x,y\) 分别对应 \(e_0,e_1\)，从而 \(\{x^{q-j}y^j\}_{j=0}^{q}\) 构成一组基。
矩阵 \(M\) 在基 \((e_0,e_1)\) 下满足 \(Me_0=e_0+e_1\)、\(Me_1=e_0\)，故其在该多项式模型中诱导代入作用
\((x,y)\mapsto (x+y,x)\)。
于是 \(x^{q-j}y^j\) 的像为 \((x+y)^{q-j}x^{j}\)，其在基 \(\{x^{q-i}y^i\}\) 下的展开系数恰为 \(\binom{q-j}{i}\)；
从而 \(\mathrm{Sym}^q(M)\) 的矩阵系数与命题 \ref{prop:pom-disjointness-symmetric-compression-bq} 的闭式一致，得到 \(B_q\cong\mathrm{Sym}^q(M)\)。
\par
由于 \(M\) 的特征值为 \(\varphi\) 与 \(-\varphi^{-1}\)，对称幂的谱为所有乘积 \(\varphi^{q-k}(-\varphi^{-1})^{k}=(-1)^k\varphi^{q-2k}\ (0\le k\le q)\)，且次数为 \(q+1\)，故重数均为 \(1\)。证毕。
\end{proof}

\begin{proposition}[两条本征线性形式生成全谱：\texorpdfstring{$B_q$}{Bq} 的闭式特征向量]\label{prop:pom-bq-explicit-diagonalization-linear-forms}
\leanverified{paper\_pom\_bq\_explicit\_diagonalization\_linear\_forms}
令 \(\varphi=(1+\sqrt5)/2\)、\(\psi=-\varphi^{-1}\)，并定义线性形式
\[
\ell_+(x,y)=\varphi x+y,\qquad \ell_-(x,y)=\psi x+y.
\]
则在代入作用 \((x,y)\mapsto(x+y,x)\) 下有
\(
\ell_\pm(x+y,x)=\lambda_\pm\,\ell_\pm(x,y)
\)，其中 \(\lambda_+=\varphi\)、\(\lambda_-=\psi\)。
对每个 \(k\in\{0,1,\dots,q\}\)，令
\[
P_{q,k}(x,y):=\ell_+(x,y)^{\,q-k}\,\ell_-(x,y)^{\,k}\in\QQ(\sqrt5)[x,y]_q.
\]
则 \(P_{q,k}\) 为 \(B_q\) 的本征向量，并满足严格本征关系
\[
\boxed{\;P_{q,k}(x+y,x)=(-1)^k\varphi^{\,q-2k}\,P_{q,k}(x,y)\;}
\]
其本征值与命题 \ref{prop:pom-bq-is-symq-and-spectrum} 的 \(\mu_k\) 一致。
在基 \(e_j=x^{q-j}y^j\) 下写 \(P_{q,k}=\sum_{j=0}^{q} v^{(k)}_j\,e_j\)，则坐标分量具有闭式
\[
\boxed{\;
v^{(k)}_j
=\sum_{a=0}^{j}
\binom{q-k}{a}\binom{k}{j-a}\,
\varphi^{\,q-k-a}\,
\psi^{\,k-j+a}\qquad(0\le j\le q)\; } .
\]
并且由于命题 \ref{prop:pom-bq-weighted-selfadjoint-binomial} 的 \(W_q\)-自伴随性与谱的单重性，向量族 \(\{v^{(k)}\}_{k=0}^{q}\) 可取为 \(\langle\cdot,\cdot\rangle_{W_q}\) 下两两正交的一组本征基。
\end{proposition}
\begin{proof}
由定义，
\[
\ell_\pm(x+y,x)=\lambda_\pm(x+y)+x=(\lambda_\pm+1)x+\lambda_\pm y=\lambda_\pm^2 x+\lambda_\pm y=\lambda_\pm(\lambda_\pm x+y)=\lambda_\pm\,\ell_\pm(x,y),
\]
其中 \(\lambda_\pm\) 满足 \(\lambda_\pm^2=\lambda_\pm+1\)，故 \(\lambda_+=\varphi\)、\(\lambda_-=\psi\)。
于是
\[
P_{q,k}(x+y,x)=\lambda_+^{\,q-k}\lambda_-^{\,k}\,P_{q,k}(x,y)=\varphi^{\,q-k}\psi^{\,k}\,P_{q,k}(x,y)=(-1)^k\varphi^{\,q-2k}P_{q,k}(x,y),
\]
给出本征关系与本征值。
再将
\(
(\varphi x+y)^{q-k}(\psi x+y)^k
\)
作二项式展开并按 \(y\) 的幂次收集，即得到 \(v^{(k)}_j\) 的所示闭式。
最后，由命题 \ref{prop:pom-bq-weighted-selfadjoint-binomial}，\(B_q\) 在 \(\langle\cdot,\cdot\rangle_{W_q}\) 下自伴随；由于本征值两两不同（命题 \ref{prop:pom-bq-is-symq-and-spectrum}），不同本征值对应的本征向量必两两正交，从而可取为 \(W_q\)-正交本征基。证毕。
\end{proof}

\begin{theorem}[Lucas 平方谱分解：\texorpdfstring{$\det(I-uB_q^2)$}{det(I-uBq^2)} 的完全因式分解]\label{thm:pom-bq-lucas-square-spectrum-factorization}
\leanverified{paper\_pom\_bq\_lucas\_square\_spectrum\_factorization}
对任意整数 \(q\ge 1\) 有严格恒等式
\[
\boxed{\;
\det(I-uB_q^2)
=(1-u)^{\mathbf 1_{\{q\ \mathrm{even}\}}}
\prod_{\substack{1\le d\le q\\ d\equiv q\ (\mathrm{mod}\ 2)}}
\bigl(u^2-L_{2d}\,u+1\bigr)\; } .
\]
特别地，当 \(q=4\) 时
\[
\det(I-uB_4^2)=(1-u)(u^2-7u+1)(u^2-47u+1),
\]
其中 \(7=L_4\)、\(47=L_8\) 给出四阶平方谱合迹的整系数来源。
\end{theorem}
\begin{proof}
由命题 \ref{prop:pom-bq-is-symq-and-spectrum}，\(B_q\) 的特征值为 \(\mu_k=(-1)^k\varphi^{q-2k}\)，故 \(B_q^2\) 的特征值为 \(\lambda_k:=\mu_k^2=\varphi^{2(q-2k)}\)。
当 \(q\) 为偶数时存在唯一自互反模态 \(k=q/2\) 使 \(\lambda_{q/2}=1\)，贡献因子 \(1-u\)；其余特征值按互反配对 \(\lambda\leftrightarrow 1/\lambda\)。
对任意正整数 \(d\equiv q\ (\mathrm{mod}\ 2)\)（对应 \(\lambda=\varphi^{2d}\)），有 Lucas 偶指标恒等式
\(
L_{2d}=\varphi^{2d}+\varphi^{-2d}
\)；
从而
\[
(1-u\varphi^{2d})(1-u\varphi^{-2d})=u^2-(\varphi^{2d}+\varphi^{-2d})u+1=u^2-L_{2d}u+1.
\]
把所有互反对的因子相乘即得结论。证毕。
\end{proof}

\begin{theorem}[闭式迹公式：\texorpdfstring{$\mathrm{Tr}(B_q^n)$}{Tr(Bq^n)} 的 Lucas 展开]\label{thm:pom-bq-trace-lucas}
\leanverified{paper\_pom\_bq\_trace\_lucas}
对任意整数 \(q\ge 1\)、\(n\ge 1\)，成立严格闭式
\[
\boxed{\;
\mathrm{Tr}(B_q^{\,n})
=
\sum_{k=0}^{\lfloor q/2\rfloor-\mathbf 1_{\{q\ \mathrm{even}\}}}
(-1)^{kn}\,L_{n(q-2k)}
\;+\;
\mathbf 1_{\{q\ \mathrm{even}\}}\cdot (-1)^{(q/2)n}\; } .
\]
因此对 \(B_q\) 的 primitive 轨道数（Möbius 反演）
\[
p_n(B_q):=\frac{1}{n}\sum_{d\mid n}\mu(d)\,\mathrm{Tr}\!\bigl(B_q^{\,n/d}\bigr)
\]
可在有限步内化为 Lucas 数的有限线性组合。
\end{theorem}
\begin{proof}
由命题 \ref{prop:pom-bq-is-symq-and-spectrum}，
\(
\mathrm{Tr}(B_q^{\,n})=\sum_{k=0}^{q}(-1)^{kn}\varphi^{n(q-2k)}
\)。
对 \(k\) 与 \(q-k\) 配对，并注意 \(q-2k\equiv q\ (\mathrm{mod}\ 2)\) 使得 \((-1)^{n(q-2k)}=(-1)^{nq}\)，可得每对之和为
\[
(-1)^{kn}\Bigl(\varphi^{n(q-2k)}+(-1)^{n(q-2k)}\varphi^{-n(q-2k)}\Bigr)=(-1)^{kn}L_{n(q-2k)}.
\]
若 \(q\) 为偶数，还需加上中点模态 \(k=q/2\) 的贡献 \((-1)^{(q/2)n}\)。
primitive 轨道数的 Lucas 表达由定义的 Möbius 反演直接代入得到。证毕。
\end{proof}

\paragraph{整数阶压力的离散凸性与二阶碰撞的协议下界}
\begin{corollary}[离散凸性与 \texorpdfstring{$\rho(A_2)$}{rho(A2)} 的严格协议下界]\label{cor:pom-integer-pressure-discrete-convexity-strict-bound}
记
\[
\tau_q:=\lim_{m\to\infty}\frac{1}{m}\log S_q(m)
\qquad(q\in\NN),
\]
其中 \(\tau_0=\log\varphi\)（因 \(|X_m|=S_0(m)\asymp \varphi^m\)）且 \(\tau_1=\log 2\)（因 \(S_1(m)=\sum_x d_m(x)=|\Omega_m|=2^m\)）。
则整数列 \((\tau_q)_{q\ge 0}\) 满足离散凸性
\[
\boxed{\;
\tau_{q+1}-\tau_q\ \ge\ \tau_q-\tau_{q-1}\qquad(q\ge 1)\; }.
\]
若进一步处于定理 \ref{thm:pom-real-q-pressure-analytic-interpolation} 的严格凸制度（例如纤维多重度势不与常数同调），则上述不等式对一切 \(q\ge 1\) 严格成立。
特别地在非平凡折叠（严格凸）情形下，
\[
\boxed{\;
\tau_2\ >\ 2\tau_1-\tau_0\ =\ 2\log 2-\log\varphi,
\qquad\text{等价地}\qquad
\rho\!\left(A_2^{\mathrm{ten}}\right)\ >\ \frac{4}{\varphi}\; }.
\]
\end{corollary}
\leanverified{paper\_pom\_integer\_pressure\_discrete\_convexity\_strict\_bound}
\begin{proof}
离散凸性为推论 \ref{cor:pom-crossq-pressure-convexity} 的整点形式（亦可直接对每个固定 \(m\) 用 Hölder 不等式得到 \(S_q(m)^2\le S_{q-1}(m)S_{q+1}(m)\)，再取对数并除以 \(m\) 取极限）。
严格性由定理 \ref{thm:pom-real-q-pressure-analytic-interpolation} 的严格凸条件推出。
将 \(q=1\) 代入即得所示二阶下界与谱半径不等式。证毕。
\end{proof}

\paragraph{二阶碰撞谱的可去因子分解：同步尾指数的谱反演}
\begin{proposition}[同步吸收块与歧义尾核：碰撞 \texorpdfstring{$\zeta$}{zeta} 的可去因子]\label{prop:pom-collision-zeta-removable-sync-factor}
\leanverified{paper\_pom\_collision\_zeta\_removable\_sync\_factor}
在命题 \ref{prop:pom-momq-tensor-kernel-construction} 的 \(q=2\) 情形，按逆码歧义态空间 \(\mathcal S\) 的“已同步态/未同步态”（候选集为单元 vs.\ 非单元）作分解，
可将 \(A_2^{\mathrm{ten}}\) 通过对基的置换重排为块三角矩阵
\[
A_2^{\mathrm{ten}}\ \sim\
\begin{pmatrix}
D & *\\
0 & B
\end{pmatrix},
\]
其中 \(D\) 为同步吸收块的计数核，\(B\) 为未同步歧义子系统的尾谱核。
于是有行列式分解
\[
\boxed{\;
\det(I-zA_2^{\mathrm{ten}})
=
\det(I-zD)\,\det(I-zB)\; }.
\]
因此，一旦同步吸收因子 \(\det(I-zD)\) 可显式计算，则“同步尾”主导极点半径 \(z_{\mathrm{sync}}:=\rho(B)^{-1}\) 可由二阶碰撞谱
\(\det(I-zA_2^{\mathrm{ten}})\)
在剔除 \(\det(I-zD)\) 后的剩余极点集合中单独反演得到。
\end{proposition}
\begin{proof}
由“已同步态”在确定性逆码更新下保持为已同步态（候选集一旦为单元即永不再分裂），同步态张成的子空间为不变子空间，从而存在如上块三角分解。
块三角矩阵的行列式等于对角块行列式之积，故得分解式；极点反演结论随之成立。证毕。
\end{proof}

\paragraph{黄金两态骨架诱导的二次特征方程：\texorpdfstring{$\rho(A_q)$}{rho(Aq)} 的固定点表征}
\begin{proposition}[Perron 根的固定点方程与单调求根]\label{prop:pom-golden-grammar-fixed-point-rho-aq}
\leanverified{paper\_pom\_golden\_grammar\_fixed\_point\_rho\_aq}
在命题 \ref{prop:pom-momq-tensor-kernel-construction} 的设定下，\(A_q^{\mathrm{ten}}\) 可写为 \(2\times 2\) 块矩阵
\[
A_q^{\mathrm{ten}}
\;=\;
\begin{pmatrix}
T_0^{\otimes q} & T_1^{\otimes q}\\[2pt]
T_0^{\otimes q} & 0
\end{pmatrix}.
\]
令 \(\lambda_q:=\rho(A_q^{\mathrm{ten}})>0\)。
则 \(\lambda_q\) 等价于唯一满足
\[
\boxed{\;
\rho\!\left(\lambda_q^{-1}T_0^{\otimes q}+\lambda_q^{-2}T_1^{\otimes q}T_0^{\otimes q}\right)=1\; }
\]
的正实数。
并且函数
\[
F_q(\lambda):=\rho\!\left(\lambda^{-1}T_0^{\otimes q}+\lambda^{-2}T_1^{\otimes q}T_0^{\otimes q}\right)
\]
在 \(\lambda>0\) 上严格单调递减，从而可用二分法或单调迭代在一维上收敛求得 \(\lambda_q\)。
\end{proposition}
\begin{proof}
对 Perron 本征对 \((\lambda_q,(u_0,u_1))\)（其中 \(u_0,u_1\ge 0\) 非零）写出本征方程：
\[
\lambda_q u_0=T_0^{\otimes q}u_0+T_1^{\otimes q}u_1,
\qquad
\lambda_q u_1=T_0^{\otimes q}u_0.
\]
消去 \(u_1=\lambda_q^{-1}T_0^{\otimes q}u_0\) 得
\(
u_0=\lambda_q^{-1}T_0^{\otimes q}u_0+\lambda_q^{-2}T_1^{\otimes q}T_0^{\otimes q}u_0
\)，
故 \(1\) 为矩阵 \(\lambda_q^{-1}T_0^{\otimes q}+\lambda_q^{-2}T_1^{\otimes q}T_0^{\otimes q}\) 的特征值；由非负性与不可约性，其谱半径必须为 \(1\)。
反向方向同理：若存在 \(\lambda>0\) 使该谱半径为 \(1\)，取 Perron 向量 \(u_0\) 并令 \(u_1=\lambda^{-1}T_0^{\otimes q}u_0\) 即构造出 \(A_q^{\mathrm{ten}}\) 的 Perron 本征向量，从而 \(\lambda=\rho(A_q^{\mathrm{ten}})\)。
\par
单调性来自谱半径对非负矩阵条目的单调性：\(\lambda\) 增大时两项 \(\lambda^{-1},\lambda^{-2}\) 同时严格减小，故 \(F_q\) 严格递减。证毕。
\end{proof}

\paragraph{Rényi 熵率的谱公式与碰撞指数的离散凹性}
\begin{corollary}[碰撞指数的凹性与 Rényi 熵率谱公式]\label{cor:pom-renyi-rate-spectral-formula-discrete-concavity}
\leanverified{paper\_pom\_renyi\_rate\_spectral\_formula\_discrete\_concavity}
对整数 \(q\ge 2\) 定义未分离尾
\[
\Pi_q(m):=\sum_{x\in X_m}w_m(x)^q=\frac{S_q(m)}{2^{mq}},
\]
并令其指数衰减率（碰撞指数）为
\[
E_q:=\lim_{m\to\infty}-\frac{1}{m}\log \Pi_q(m).
\]
若 \(S_q(m)=\Theta(\rho(A_q^{\mathrm{ten}})^m)\)（例如 \(A_q^{\mathrm{ten}}\) 原始且边界向量非退化），则
\[
\boxed{\;
E_q=q\log 2-\log\rho(A_q^{\mathrm{ten}})\;=\;h_q\;=\;-P(q)\; }
\]
（记号见定义 \ref{def:pom-moment-pressure}）。
并且整数列 \((E_q)_{q\ge 2}\) 在 \(q\) 上满足离散凹性
\[
\boxed{\;
E_{q+1}-E_q\ \le\ E_q-E_{q-1}\qquad(q\ge 3)\; }.
\]
相应的 Rényi 熵率（标准归一化，记为 \(h_q^{\mathrm R}\)）满足谱闭式
\[
\boxed{\;
h_q^{\mathrm R}
:=\lim_{m\to\infty}\frac{1}{m}\cdot\frac{1}{1-q}\log\sum_{x\in X_m}w_m(x)^q
\;=\;\frac{E_q}{q-1}
\;=\;\frac{q\log 2-\log\rho(A_q^{\mathrm{ten}})}{q-1}\; }.
\]
\end{corollary}
\begin{proof}
谱闭式由 \(\Pi_q(m)=S_q(m)/2^{mq}\) 与 \(\tau(q)=\log\rho(A_q^{\mathrm{ten}})\) 直接得到，并与 \(h_q=\log(2^q/r_q)\) 的定义一致（定义 \ref{def:pom-moment-pressure}）。
离散凹性由 \(\tau_q\) 的离散凸性（推论 \ref{cor:pom-integer-pressure-discrete-convexity-strict-bound} 的非严格版本）经恒等变换 \(E_q=q\log2-\tau_q\) 立即推出。证毕。
\end{proof}
